%%% ==============================================================================
%%%  ПРЕАМБУЛА СРД (по образцу отчёта по учебной практике ИТМО)
%%%  Сборка: XeLaTeX. Шрифт: Times New Roman (системный).
%%% ==============================================================================

%%% ---- Шрифт: Times New Roman (XeLaTeX + fontspec) ----
\usepackage{fontspec}
\IfFontExistsTF{Times New Roman}{%
  \setmainfont{Times New Roman}[Ligatures=TeX, Numbers=OldStyle]
  \setmonofont{Times New Roman}[Scale=MatchLowercase]
}{%
  \setmainfont{TeX Gyre Termes}[Ligatures=TeX, Numbers=OldStyle]
  \setmonofont{TeX Gyre Termes}[Scale=MatchLowercase]
}

%%% ---- Язык и переносы ----
\usepackage[english,russian]{babel}

%%% ---- Математика ----
\usepackage{amsmath}

%%% ---- Геометрия страницы (ГОСТ, методические указания ИТМО) ----
\usepackage{geometry}
\geometry{
  a4paper,
  left   = 30mm,
  right  = 15mm,
  top    = 20mm,
  bottom = 20mm,
}

%%% ---- Межстрочный интервал: полуторный (1.5) ----
\usepackage{setspace}
\onehalfspacing

%%% ---- Абзацный отступ 1,25 см ----
\usepackage{indentfirst}
\setlength{\parindent}{1.25cm}

\sloppy

%%% ---- Нумерация страниц ----
\usepackage{fancyhdr}
\pagestyle{fancy}
\fancyhf{}
\fancyfoot[C]{\thepage}
\renewcommand{\headrulewidth}{0pt}
\renewcommand{\footrulewidth}{0pt}
\fancypagestyle{plain}{%
  \fancyhf{}%
  \fancyfoot[C]{\thepage}%
  \renewcommand{\headrulewidth}{0pt}%
  \renewcommand{\footrulewidth}{0pt}%
}

\setcounter{secnumdepth}{4}
\setcounter{tocdepth}{3}

%%% ---- Заголовки разделов (titlesec) ----
\usepackage{titlesec}
\newcommand{\nohyphen}{%
  \tolerance=10000\hyphenpenalty=10000\exhyphenpenalty=10000%
}

\titleformat{\section}
  {\clearpage\nohyphen\bfseries\normalsize}
  {\thesection}
  {0.5em}
  {}
\titlespacing*{\section}
  {\parindent}{*2}{*1}

\titleformat{\subsection}
  {\nohyphen\bfseries\normalsize}
  {\thesubsection}
  {0.5em}
  {}
\titlespacing*{\subsection}
  {\parindent}{*2}{*1}

\titleformat{\subsubsection}
  {\nohyphen\bfseries\normalsize}
  {\thesubsubsection}
  {0.5em}
  {}
\titlespacing*{\subsubsection}
  {\parindent}{*2}{*1}

\titleformat{\paragraph}
  {\nohyphen\bfseries\normalsize}
  {\theparagraph}
  {0.5em}
  {}
\titlespacing*{\paragraph}
  {\parindent}{*1}{*0.5}

%%% ---- Содержание (tocloft) ----
\usepackage{tocloft}
\addto\captionsrussian{\renewcommand{\contentsname}{СОДЕРЖАНИЕ}}
\renewcommand{\cftsecleader}{\cftdotfill{\cftdotsep}}
\setlength{\cftsecindent}{0em}
\setlength{\cftsecnumwidth}{2em}
\setlength{\cftsubsecindent}{1.25cm}
\setlength{\cftsubsecnumwidth}{3em}
\setlength{\cftsubsubsecindent}{2.50cm}
\setlength{\cftsubsubsecnumwidth}{4em}
\setlength{\cftbeforesecskip}{0pt}
\renewcommand{\cftsecfont}{\normalfont}
\renewcommand{\cftsecpagefont}{\normalfont}

\makeatletter
\renewcommand\tableofcontents{%
  \clearpage
  \phantomsection
  {\centering\bfseries\contentsname\par}%
  \vspace{1em}%
  \@starttoc{toc}%
}
\makeatother

%%% ---- Подписи рисунков и таблиц ----
\usepackage{caption}
\DeclareCaptionLabelSeparator{dash}{~--- }
\captionsetup[figure]{
  name = {Рисунок},
  labelsep = dash,
  justification = centering,
  font = {singlespacing},
  position = below,
}
\captionsetup[table]{
  name = {Таблица},
  labelsep = dash,
  justification = centering,
  font = {singlespacing},
  position = below,
  skip = 0.5em,
}

\numberwithin{equation}{section}
\counterwithin{figure}{section}
\counterwithin{table}{section}

%%% ---- Графика ----
\usepackage{graphicx}
\usepackage{float}
\graphicspath{{img/},{../}}

%%% ---- Таблицы ----
\usepackage{tabularx}
\usepackage{longtable}
\usepackage{booktabs}
\usepackage{multirow}
\usepackage{array}

%%% ---- Списки ----
\usepackage{enumitem}
\setlist{nosep, leftmargin=\parindent}
\setlist[itemize]{label=--}
\setlist[enumerate]{leftmargin=\parindent, labelsep=0.5em}

%%% ---- Обход отсутствующего шрифта pzdr (hyperref в XeLaTeX) ----
\makeatletter
\let\srd@orig@font\font
\def\font{\srd@font@check}
\def\srd@font@check#1=#2 #3 #4{%
  \ifx#1\XeTeXLink@font
    \ifx#2pzdr
      \srd@orig@font\XeTeXLink@font=cmr10 at 1sp\relax
    \else
      \srd@orig@font#1=#2 #3 #4
    \fi
  \else
    \srd@orig@font#1=#2 #3 #4
  \fi
}
\makeatother

%%% ---- Гиперссылки ----
\usepackage{hyperref}
\hypersetup{
  colorlinks = true,
  linkcolor  = black,
  citecolor  = black,
  urlcolor   = blue,
  hidelinks,
}

%%% ---- Восстановление примитива \font после hyperref ----
\makeatletter
\let\font\srd@orig@font
\makeatother

\newcommand{\figref}[1]{рисунок~\ref{#1}}
\newcommand{\tabref}[1]{таблица~\ref{#1}}

\fancypagestyle{titlepage}{%
  \fancyhf{}%
  \renewcommand{\headrulewidth}{0pt}%
  \renewcommand{\footrulewidth}{0pt}%
}
